%Daten zur Institution

%\input{Dozentdaten}

%\renewcommand{\fachbereich}{Fachbereich}

%\renewcommand{\dozent}{Prof. Dr. . }

%Klausurdaten

\renewcommand{\klausurgebiet}{ }

\renewcommand{\klausurtyp}{ }

\renewcommand{\klausurdatum}{ . 20}

\klausurvorspann {\fachbereich} {\klausurdatum} {\dozent} {\klausurgebiet} {\klausurtyp}

%Daten für folgende Punktetabelle

\renewcommand{\aeins}{  3  }

\renewcommand{\azwei}{  3   }

\renewcommand{\adrei}{  7  }

\renewcommand{\avier}{  4  }

\renewcommand{\afuenf}{  0  }

\renewcommand{\asechs}{  0  }

\renewcommand{\asieben}{  0  }

\renewcommand{\aacht}{  5  }

\renewcommand{\aneun}{  15  }

\renewcommand{\azehn}{  5  }

\renewcommand{\aelf}{  0  }

\renewcommand{\azwoelf}{  3  }

\renewcommand{\adreizehn}{  3   }

\renewcommand{\avierzehn}{  0  }

\renewcommand{\afuenfzehn}{  10  }

\renewcommand{\asechzehn}{  4  }

\renewcommand{\asiebzehn}{  62  }

\renewcommand{\aachtzehn}{    }

\renewcommand{\aneunzehn}{    }

\renewcommand{\azwanzig}{    }

\renewcommand{\aeinundzwanzig}{    }

\renewcommand{\azweiundzwanzig}{    }

\renewcommand{\adreiundzwanzig}{    }

\renewcommand{\avierundzwanzig}{    }

\renewcommand{\afuenfundzwanzig}{    }

\renewcommand{\asechsundzwanzig}{    }

\punktetabellesechzehn

\klausurnote

\newpage

\setcounter{section}{K}

\inputaufgabegibtloesung
{3}
{

Definiere die folgenden
\zusatzklammer {kursiv gedruckten} {} {} Begriffe.
\aufzaehlungsechs{Eine
\stichwort {bogenparametrisierte} {}
Kurve
\maabbdisp {\gamma} {[a,b]} { \R^n
} {.}

}{Zwei in einem Punkt

\mavergleichskette
{\vergleichskette
{  P
}
{ \in }{ M
}
{  }{}
{    }{}
{   }{}
}
{}{}{}
einer differenzierbaren Mannigfaltigkeit $M$ \stichwort {tangential äquivalente} {} differenzierbare Kurven
\maabbdisp {\gamma_1, \gamma_2} { I } { M
} {}
\zusatzklammer {dabei sei

\mavergleichskettek
{\vergleichskettek
{  0
}
{ \in }{  I
}
{  }{
}
{    }{
}
{   }{
}
}
{}{}{}
ein offenes reelles Intervall und

\mavergleichskettek
{\vergleichskettek
{   \gamma_1(0)
}
{ = }{ \gamma_2(0)
}
{ = }{ P
}
{    }{
}
{   }{
}
}
{}{}{}} {} {.}

}{Ein
\stichwort {stetiger Schnitt} {}
zu einer stetigen Abbildung
\maabbdisp {p} {Y} {X
} {}
zwischen topologischen Räumen
\mathkor {} {X} {und} {Y} {.}

}{Die \stichwort {kanonische Volumenform} {} auf einer orientierten riemannschen Mannigfaltigkeit $M$.

}{Die
\stichwort {äußere Ableitung} {}
einer
\definitionsverweis {stetig differenzierbaren}{}{}
$k$-\definitionsverweis {Differentialform}{}{}

\mathl{\omega \in
{ \mathcal E }^{ k } (  U   )}{} auf einer
\definitionsverweis {offenen Menge}{}{}

\mathl{U \subseteq \R^n}{.}

}{Die \stichwort {Schnittkrümmung} {} zu linear unabhängigen Vektoren

\mavergleichskette
{\vergleichskette
{  v,w
}
{ \in }{T_PM
}
{  }{
}
{    }{
}
{   }{
}
}
{}{}{}
auf einer
\definitionsverweis {riemannschen Mannigfaltigkeit}{}{}
$M$ in einem Punkt

\mavergleichskette
{\vergleichskette
{ P
}
{ \in }{ M
}
{  }{
}
{    }{
}
{   }{
}
}
{}{}{.}
}

}
{} {}

\inputaufgabegibtloesung
{3}
{

Formuliere die folgenden Sätze.
\aufzaehlungdrei{Der
\stichwort {Satz über die Orientierungen auf einer differenzierbaren Hyperfläche} {}

\mavergleichskette
{\vergleichskette
{  Y
}
{ \subseteq }{  \R^n
}
{  }{
}
{    }{
}
{   }{
}
}
{}{}{.}}{Die Formel für die Lie-Klammer von Vektorfeldern auf einer offenen Menge

\mavergleichskette
{\vergleichskette
{  U
}
{ \subseteq }{  \R^n
}
{  }{
}
{    }{
}
{   }{
}
}
{}{}{.}}{Der Satz über
\stichwort {Retraktionen zum Rand} {}
auf Mannigfaltigkeiten.}

}
{} {}

\inputaufgabegibtloesung
{7}
{

Beweise den
\stichwort {Satz über die Lösung von gewöhnlichen Differentialgleichungen auf einer differenzierbaren Hyperfläche} {}

\mavergleichskette
{\vergleichskette
{ Y
}
{ \subseteq }{ \R^n
}
{  }{
}
{    }{
}
{   }{
}
}
{}{}{.}

}
{} {}

\inputaufgabegibtloesung
{4}
{

Es sei
\maabb {h} { I } { V
} {}
eine zweimal differenzierbare Kurve in einem
\definitionsverweis {euklidischen Vektorraum}{}{}
$V$. Zeige, dass bei

\mavergleichskette
{\vergleichskette
{h'(t)
}
{ \neq }{ 0
}
{  }{
}
{    }{
}
{   }{
}
}
{}{}{}
die Gleichheit

\mavergleichskettedisp
{\vergleichskette
{ \Vert { h'(t)} \Vert'
}
{ =} { { \frac{   \left\langle h'(t) , h^{\prime \prime} (t)  \right\rangle  }{  \left\langle h'(t) , h'(t)  \right\rangle  } } \Vert { h'(t)} \Vert
}
{ } {
}
{ } {
}
{ } {
}
}
{}{}{}
gilt.

}
{} {}

\inputaufgabe
{0}
{

}
{} {}

\inputaufgabe
{0}
{

}
{} {}

\inputaufgabe
{0}
{

}
{} {}

\inputaufgabegibtloesung
{5}
{

Beweise die Aussage, dass die tangentiale Äquivalenz von Wegen auf einer Mannigfaltigkeit in einem Punkt $P$ mit einer beliebigen Karte überprüft werden kann.

}
{} {}

\inputaufgabegibtloesung
{15 (3+3+4+5)}
{

Es sei

\mavergleichskette
{\vergleichskette
{T
}
{ = }{S^1 \times S^1
}
{  }{
}
{    }{
}
{   }{
}
}
{}{}{}
der Torus und

\mavergleichskette
{\vergleichskette
{S^2
}
{ \subset }{\R^3
}
{  }{
}
{    }{
}
{   }{
}
}
{}{}{}
die
\definitionsverweis {Einheitssphäre}{}{.}

a) Zeige, dass durch
\maabbeledisp {\varphi} {S^1 \times S^1} { S^2
} { \begin{pmatrix}  \alpha \\ \beta   \end{pmatrix} } { { \frac{   1 }{  2 } } \begin{pmatrix}  \sqrt{ 2-2 \sin \alpha   } \cos \beta  \\ \cos \alpha + \sin \beta \cos \alpha \\  1  + \sin \alpha -  \sin \beta + \sin \alpha \sin \beta   \end{pmatrix}
} {}
eine stetige Abbildung gegeben ist.

b) Zeige, dass $\varphi$ surjektiv ist.

c) Beschreibe die Fasern von $\varphi$.

d) Erläutere die Abbildung $\varphi$ unter Verwendung einer Skizze.

}
{} {}

\inputaufgabegibtloesung
{5}
{

Berechne das Wegintegral
\mathl{\int_\gamma \omega}{} zu
\maabbeledisp {\gamma} { [-1,0] } {\R^3
} { t } {(-t^2,t^3-1,t+2)
} {,}
für die $1$-Differentialform

\mavergleichskettedisp
{\vergleichskette
{ \omega
}
{ =} { x^3dx -yzdy +xz^2 dz
}
{ } {
}
{ } {
}
{ } {
}
}
{}{}{}
auf dem $\R^3$.

}
{} {}

\inputaufgabe
{0}
{

}
{} {}

\inputaufgabegibtloesung
{3 (2+1)}
{

Es sei
\maabbeledisp {f} {\R} {\R
} { t } {f(t)
} {,}
die durch

\mavergleichskettedisp
{\vergleichskette
{ f(t)
}
{ =} { { \frac{   \sin^{ 3  } (t^4)  }{  1+t^2 } }
}
{ } {
}
{ } {
}
{ } {
}
}
{}{}{}
gegeben ist.

a) Berechne die äußere Ableitung von $f$.

b) Berechne die äußere Ableitung von $fdt$.

}
{} {}

\inputaufgabegibtloesung
{3}
{

Man gebe ein Beispiel für eine stetig differenzierbare Abbildung
\maabbdisp {\varphi} { H } { H
} {}
einer euklidischen Halbebene in sich mit der Eigenschaft, dass genau ein Randpunkt von $H$ in einen Randpunkt und alle anderen Punkte in das Innere der Halbebene abgebildet werden.

}
{} {}

\inputaufgabe
{0}
{

}
{} {}

\inputaufgabegibtloesung
{10}
{

Beweise die Quaderversion des Satzes von Stokes.

}
{} {}

\inputaufgabegibtloesung
{4}
{

Bestimme für die Kurve
\maabbeledisp {\psi} {\R} { {\mathbb H}
} { t } { \left(  t   , \,  e^t                   \right)
} {,}
in die
\definitionsverweis {Halbebene von Poincaré}{}{}
${\mathbb H}$ die
\definitionsverweis {Christoffelsymbole}{}{}
für den
\definitionsverweis {zurückgezogenen Zusammenhang}{}{}
\zusatzklammer {zum
\definitionsverweis {Levi-Civita-Zusammenhang}{}{}
auf ${\mathbb H}$} {} {}
auf $\R$.

}
{} {}

 [[Kategorie:Latexseite]]
